% !PRINTED 295
\weberpage{295}
%<<CONT
Möglichkeit gegeben, aus den Vorzeichen der Functionen $\varphi, \psi$
in den Theilpunkten der Strecken die Charakteristik vollständig
zu bestimmen. Wie wir aber die Curve $f$ in solche Theilstrecken
eintheilen, das lehrt uns der \gesperrt{Sturm'sche Lehrsatz}. Wir denken
uns zu diesem Zweck $x$ und $y$ als Functionen einer Variablen $t$ dargestellt, die längs der Curve $f$ alle
reellen Werthe von $-\infty$ bis $+\infty$
(oder auch nur die Werthe eines
endlichen Intervalles) durchläuft. Dadurch gehen $\varphi$ und $\psi$ auch in Functionen der Variablen $t$ über und man
hat die Lage ihrer Wurzeln nur nach dem \gesperrt{Sturm'schen Lehrsatz} zu untersuchen.

\weberfigureimage{weber-b1-fig14.png}{0.44}

Ist z.\,B. wie in Fig. 14 die Curve $f$ ein Kreis mit der
Gleichung
\[
(x - \alpha)^2 + (y - \beta)^2 = \varrho^2,
\]
so können wir setzen:
\[
x - \alpha = \varrho\, \frac{1 - t^2}{1 + t^2},
\]
\[
y - \beta = \varrho\, \frac{2t}{1 + t^2},
\]
und $t$ durchläuft einfach alle Werthe von $-\infty$ bis $+\infty$, während
der Punkt $x, y$ sich über den ganzen Kreis in positivem Sinne
bewegt. $\varphi$ und $\psi$ gehen, mit einer geeigneten Potenz von
$1 + t^2$ multiplicirt, in ganze rationale Functionen von $t$ über,
die die Anwendung des \gesperrt{Sturm'schen Satzes} gestatten.

\begin{center}
\S.\,98.

\gesperrt{Gauss' erster Beweis des Fundamentalsatzes der Algebra.}
\end{center}

\gesperrt{Gauss} hat für den Hauptsatz der Algebra, dass eine Gleichung
$n^{\text{ten}}$ Grades $n$ Wurzeln hat, drei verschiedene Beweise
gegeben. Der erste von diesen, der sich in seiner Doctordissertation findet, und den er später noch weiter ausgeführt und
anders dargestellt hat, beruht, wenn auch in anderer Einkleidung,
%CONT>>
